\Needspace{6\baselineskip}\pdfbookmark[3]{III. Valuation of obligations towards future members}{witwen-II-2-III}
\par\medskip\noindent III.\quad Valuation of the present monetary value of the fund's obligations towards all members joining in the future.\par\smallskip
For this third heading, the fund's past experience offers a highly valuable aid; indeed, without this experience an estimate would be quite impossible. My procedure is as follows.

For each member who joined by Michaelmas 1804, up to and including No. 108\footnote{My decision to go precisely this far and no further was deliberate, but it would take me too far to develop the reasons here in detail.}, I have reduced both the value of the contributions and that of the pensions received by their surviving dependants to the time of their entry, discounting at two different rates of interest, but with the following modifications, necessary for the use to be made of these calculations:

1) The annual contributions and the pensions have been reckoned, not as they were actually paid, but at their present levels, namely 10 Rthl. for the former and 250 Rthl. for the latter.

2) The orphans' pensions, which until 1794 were paid only until the youngest child had completed its 12\textsuperscript{th} year, have been reckoned as though they had lasted 8 years longer. Through investigat%
\SourceLeaf{159}ion in the court archives it has been established that all the youngest children concerned actually reached the age of 20. It may, however, be that somewhat too little has still been reckoned in this way, since at the death of one or other widowed professor or professor's widow there may possibly still have been children between 12 and 20 years of age, who would therefore have been entitled to pensions under the present arrangements and should thus have been included in my calculation: at least, the court records have not established the contrary with certainty in every case. Nevertheless, in any event this could not have produced a substantial increase in the total sum.

3) Among those 108 members there are 7 whose widows are still alive now. For these, I have added to the pension payments they have received so far the future amount already valued above by the calculus of probabilities, after discounting it back from the date 1845 October 1 to the time when the respective husbands joined. Without this procedure I should have had not merely to exclude these members, but, to be consistent, instead of members 1—108, to restrict myself to 1—60, a matter I cannot develop further here.

To save space, I give here not the results of the 108 calculations individually, but only their combined total. One finds
\begin{center}
\footnotesize\setlength{\tabcolsep}{3pt}
\begin{tabular}{l r|r}
&\multicolumn{2}{c}{at the interest rate}\\
&\(3\tfrac{1}{2}\) per cent&4 per cent\\\cline{2-3}
Sum of the discounted contributions . . . . .&13909,00 Thl.&13115,94 Thl.\\
—\quad —\quad —\quad Widows' pensions&67915,80 ,, &59019,67 ,, \\
—\quad —\quad —\quad Orphans' pensions&10945,04 ,, &9878,34 ,, \\
Sum of all discounted pensions . . .&78860,84 ,, &68898,01 ,,
\end{tabular}
\end{center}
\noindent\hspace*{\parindent}It follows first that the (reduced) amount of the orphans' pensions is to be taken as very nearly equal to a sixth of the (reduced) amount of the widows' pensions: in fact, at a rate of \(3\tfrac{1}{2}\) per cent the former is somewhat smaller, but at a rate of 4 per cent somewhat larger, than this \(\tfrac{1}{6}\). It could easily be shown that this small variation in the result with the rate of interest is entirely grounded in the nature of the matter. Since, however, the body of experience is too small for minute precision to be ascribed to the result at all, I have, in the calculation above on p. [152], simply kept to the simple ratio \(1:6\).

One may further say that the combined rights of all those 108 members, valued at the present levels of pensions and contributions and reduced to the times of their respective entries, had the monetary value
\begin{center}
\begin{tabular}{l}
of 64951,84 thalers at an interest rate of \(3\tfrac{1}{2}\) per cent\\
\phantom{of }55782,07 thalers at an interest rate of 4 per cent
\end{tabular}
\end{center}
\noindent or, insofar as such extensive experience can serve as a guide, that on average entry into the fund has a monetary value of 601,40 thalers, or 516,50 thalers, for each individual, according as one reckons interest at \(3\tfrac{1}{2}\) or 4 per cent.

Some readers little acquainted with matters of this kind may perhaps be surprised by the smallness of this result. But much of what seems surprising disappears when one considers that the reduction to each individual's time of entry is the basis here, and that any sum of money appears enormously reduced when discounted over a considerable number of years. I shall illustrate this by the example of \textsc{Heyne}, who paid contributions for 48 years and whose widow received a pension for 21 years and 11 months. On my bases of calculation, the sum of the contributions thus comes to 480 Rthl. and that of the pension receipts to 5479 Rthl. 4 Ggr., and hence the
\SourceLeaf{160}
\textit{net} gain to the family through admission to the fund (not as it actually was, but as it would have been at the present levels of contributions and pensions) comes to 4999 Rthl. 4 Ggr. But if those sums are reduced by discounting at 4 per cent to the time of entry, Michaelmas 1763, they fall as follows:
\begin{center}
\begin{tabular}{l l r l r l}
the contributions&to&211&Rthl.&23&Ggr.\\
the pension receipts&to&527&,,&10&,,
\end{tabular}
\end{center}
\Needspace{24\baselineskip}
\noindent and hence that great gain, in its monetary value at that date, falls to the small sum of 315 Rthl. 11 Ggr. A check of the calculation would show exactly that this small capital, increased annually by the addition of 10 Rthl. and by interest at 4 per cent, and thus continually growing, would have reached such a size at the end of 1812 that from then on the further interest and, in part, withdrawals of capital covered the pension exactly until the end of November 1834.

For this reason, as well as for others whose development would take me too far here, the mean monetary value of admission to the society, in the sense in which it is understood here, is not at all comparable with the average monetary value of the claims of all the members coexisting at a given time, among whom there is always a large proportion who have already been in the society for a very long time.

